% !TeX root = elegantnote-cn.tex
\section{胶体分散系统和大分子溶液}

\begin{enumerate}
    \item 把一种或几种物质分散在另一种物质中，构成分散系统。分散系统中被分散的物质称为\textbf{分散相}，分散分散相的物质称为\textbf{分散介质}。
    \item 按分散相粒子大小，常有三种分散系统：\textbf{分子（离子）分散系统}（$r < \qty{1}{\nano\meter}$）、\textbf{胶体分散系统}（$\qty{1}{\nano\meter} \leq r \leq \qty{100}{\nano\meter}$） 和 \textbf{粗分散系统}（$r > \qty{100}{\nano\meter}$）。
\end{enumerate}
\subsection{胶体和胶体的基本特性 20260527 20260601}

\subsubsection{分散系统的分类}

\begin{enumerate}
    \item 胶体系统包含性质不同的两大类
    \begin{enumerate}
        \item 憎液溶胶：难溶物分散在分散介质中。不稳定，易被破坏而聚沉，聚沉后往往不能恢复原态。
        \item 亲液溶胶
    \end{enumerate}
    \item 溶胶系统也可按分散相和分散介质的聚集状态分类。
    \begin{enumerate}
        \item 液体作为分散介质，液溶胶：液固、液液、液气溶胶
        \item 固体作为分散介质，固溶胶：固固、固液、固气溶胶
        \item 气体作为分散介质，气溶胶：气固、气液溶胶，没有气气溶胶（单相均一体系）
    \end{enumerate}
    \item 憎液溶胶的特性
    \begin{enumerate}
        \item 特有的分散程度：粒子大小在\qtyrange{1}{100}{\nano\meter}之间
        \item 多相不均匀性：介质间有明显的相界面
        \item 易聚结不稳定性：表面自由能高，是热力学不稳定系统，自动聚结为大粒子。
    \end{enumerate}
\end{enumerate}

\subsubsection{胶团的结构}

\begin{enumerate}
    \item 憎液溶胶的必要条件：溶解度小、稳定剂
    \item 一定量难溶物分子聚结形成胶粒中心，为胶核
    \item 胶核选择性吸附组成类似的离子，形成胶粒
    \item 胶粒与介质中的相反粒子构成胶团
    \item \ch{AgI} 胶团的构造： %pic 
    \ch{[(AgI)_{$m$} $n$ I^-, \ ($n-x$) K^+]^{$x$-} \ $x$ K^+}
    \item 不同的稳定剂形成的胶团不同。
    \item 胶粒的形状对性质有重要影响。球形流动性好；带状流动性差，易产生触变现象。
\end{enumerate}

\subsection{溶胶的制备与净化 20260601}

\subsubsection{溶胶的制备}

\begin{enumerate}
    \item 必要条件：分散相离子大小落在胶体分散系范围内，加入稳定剂
    \item 用分散法和凝聚法直接制出的粒子称为原级粒子。视具体制备条件不同，这些粒子又可以聚集成较大的次级粒子。通常所制备的溶胶中粒子的大小不均一，是多级分散系统。
    \item 分散法：用机械、化学方法使得固体粒子变小……
    \begin{enumerate}
        \item 研磨法：机械粉碎，适用于脆而易碎的物质，对于柔韧性的物质必须先硬化后再粉碎。
        \item 胶溶法：将新鲜的凝聚胶粒重新分散在介质中形成溶胶，并加入适当的稳定剂，亦称胶溶剂。
        \item 超声波分散法：只用来制备乳状液。
        \item 电弧法：主要用于制备金、银、铂等金属溶胶，包括先分散（金属蒸发）后凝聚两个过程。
        \item 气相沉积法：惰性气氛中用热源将要制备成纳米级粒子的材料气化。
    \end{enumerate}
    \item 凝聚法：使分子或离子聚结成胶粒……
    \begin{enumerate}
        \item 化学凝聚法：通过化学反应使生成物过饱和，使初生成的难溶物微粒结合成胶粒，在少量稳定剂存在下形成溶胶，稳定剂一般是某一过量反应物。
        
        例如：
        \begin{enumerate}
            \item 复分解反应制硫化砷溶胶 \\
            $\ch{As2O3} + \ch{3 H2S} \ch{->} \ch{As2S3(溶胶)} + \ch{3 H2O}$
            \item 还原反应制金溶胶 \\
            $\ch{2 HAuCl4(稀溶液)} + \ch{3 HCHO(少量)} + \ch{11 KOH} \ch{->[加热]} \ch{2 Au(溶胶)} + \ch{3 HCOOK} + \ch{8 HCl} + \ch{8 H2O}$
            \item 水解反应制氢氧化铁溶胶 \\
            $\ch{FeCl3} + \ch{3 H2O(热)} \ch{<=>} \ch{Fe(OH)_3(溶胶)} + \ch{3 HCl}$
            \item 氧化还原反应制备硫溶胶 \\
        \end{enumerate}
        \item 物理凝聚法：蒸气骤冷法
        \item 更换溶剂法：利用物质在不同溶剂中溶解度的显著差别来制备溶胶，而且两种溶剂要能完全互溶。
    \end{enumerate}
\end{enumerate}

\subsubsection{溶胶的净化}

\begin{enumerate}
    \item 过多的电解质存在会使溶胶不稳定，容易聚沉。
    \item 渗析法
    \begin{enumerate}
        \item 简单渗析：将需要净化的溶胶放半透膜制成的容器内，膜外放纯溶剂。利用浓差进行渗透。
        \item 电渗析：在装有溶胶的半透膜两侧外加电场，使多余的电解质离子向相应的电极作定向移动。
    \end{enumerate}
    \item 超过滤法：用半透膜作过滤膜，利用吸滤或加压的方法使胶粒与含有杂质的介质在压差作用下迅速分离。将半透膜上的胶粒迅速用含有稳定剂的介质再次分散。可施加电压加快过滤速度、降低超过滤压力。
\end{enumerate}

\subsubsection{溶胶的形成条件和老化机理}

\begin{enumerate}
    \item 溶胶形成的过程中经历两个阶段：晶核的形成和晶体的生长
    \item 晶核形成过程的速率决定于形成和生长两个因素
    \begin{enumerate}
        \item 从溶液中析出固体的速率，即晶核形成的速率
        \begin{equation}
            v_1 = k \frac{Q - s}{s}
        \end{equation}
        \item 晶体长大的速率
        \begin{equation}
            v_2 = D A \frac{Q - s}{\delta}
        \end{equation}
    \end{enumerate}
    \item 要得到分散度很高的溶胶，则必需控制两者的值，使 $v_2$ 很小或接近于零。
    \item 当 $\frac{Q - s}{s}$ 的值很大，有利于形成溶胶；较小，有利于生成大块沉淀；很小，有利于形成溶胶。
    \item 纯化后的胶粒也会随时间慢慢增大，最终沉淀。此过程成为溶胶的老化，是自发过程。
\end{enumerate}

\subsubsection{均分散胶体的制备和应用}

\begin{enumerate}
    \item 严格控制的条件下，有可能制备出形状相同、尺寸相差不大的沉淀颗粒，组成均分散系统。颗粒的尺寸在胶体颗粒范围之内的均分散系统则称为均分散胶体系统。
    \item Einstein 理论
    \begin{equation*}
        D = \frac{R T}{L}\frac1{6\pi \eta r}
    \end{equation*}
\end{enumerate}

\subsection{溶胶的动力性质 20260601}

\begin{enumerate}
    \item 动力性质
\end{enumerate}

\subsubsection{Brown 运动}

\begin{enumerate}
    \item 粒子粒径越小，温度越高，Brown 运动越剧烈。
    \item 认为 Brown 运动是分散介质分子以不同大小和方向的力对胶体粒子不断撞击而产生的。
    \item 当半径大于 \qty{5}{\micro\meter}，Brown 运动消失。
    \item Brown 运动的本质：\textbf{Einstein-Brown 运动公式}
    \begin{equation}
        \overline{x^2} = \frac{R T}{L} \frac{t}{3\pi \eta r}
    \end{equation}
    把粒子的位移与粒子的大小、介质黏度、温度以及观察时间等联系起来。
\end{enumerate}

\subsubsection{扩散和渗透压}

\begin{enumerate}
    \item 宏观上观察到胶体粒子从高浓度向低浓度区迁移，\textbf{扩散作用}。
    \item 稀溶液中，……，浓度梯度 $\frac{\mathrm{d}c}{\mathrm{d}x}$，扩散速度 $\frac{\mathrm{d}n}{\mathrm{d}t}$；扩散速度与浓度梯度及 $AB$ 面面积 $A$ 成正比
    \begin{equation}
        \frac{\mathrm{d}n}{\mathrm{d}t} = -D A \frac{\mathrm{d}c}{\mathrm{d}x}
    \end{equation}
    上式（）即为 \textbf{Fick 第一定律}。
    负号代表从高浓度向低浓度扩散，保证扩散的值为正值。
    \item Fick 第一定律适用于浓度梯度不变的情况。实际扩山中浓度梯度随位移发生变化。
    
    设进入 $AB$ 面的扩散量 $-D A \frac{\mathrm{d}c}{\mathrm{d}x}$，离开 $EF$ 面的扩散量 $-D A \left[\frac{\mathrm{d}c}{\mathrm{d}x} + \frac{\mathrm{d}}{\mathrm{d}x} \left(\frac{\mathrm{d}c}{\mathrm{d}x}\right) \mathrm{d}x\right]$
    \begin{equation}
        \frac{\mathrm{d}c}{\mathrm{d}t}
        = \frac{D A \left[\dfrac{\mathrm{d}}{\mathrm{d}x} \left(\dfrac{\mathrm{d}c}{\mathrm{d}x}\right) \mathrm{d}x\right]}{A \cdot \mathrm{d}x}
        = D \frac{\mathrm{d^2}c}{\mathrm{d}x^2}
    \end{equation}
\end{enumerate}